\section{Residue fields and normalized length}\label{sec:lengths}

The comparison arguments require algebras on which positive parameter
Koszul homology has length zero and parameter quotients have prescribed
positive length.  The finite approximations in Section~\ref{sec:positivity}
also require that this length be defined on underlying modules over a
flat tower, without an action of the full algebra.  We construct these
objects after arranging a perfect residue field; neither unramifiedness
nor Cohen--Macaulayness of the domain is required.
Both characteristics are needed: a prime quotient of the
mixed-characteristic ambient ring may itself have characteristic $p$,
even though ambient equicharacteristic positivity is already known.

\subsection{Extending the residue field}

\begin{lemma}[Residue-field extension]\label{len:residue-extension}
Let $(B,\mathfrak b,k)$ be a complete Noetherian local ring, and let
$\overline{k}$ be an algebraic closure of $k$.  There is a complete
Noetherian local ring $(B',\mathfrak b',\overline{k})$ and a flat local map
$B\to B'$ such that $\mathfrak bB'=\mathfrak b'$.
For every finite $B$-module $V$ and every $a\geq1$,
\[
 \length_{B'}\bigl((V\otimes_B B')/(\mathfrak b')^a
                          (V\otimes_B B')\bigr)
       =\length_B(V/\mathfrak b^aV).
\]
Consequently module dimensions are preserved, as are finite lengths.
If $B$ is regular, then $B'$ is regular.  For finite $B$-modules $M,N$
whose tensor product has finite length, every Tor length, and hence
their intersection multiplicity when defined, is preserved.
\end{lemma}

\begin{proof}
In equal characteristic, a coefficient field gives a presentation
$B=k[[X_1,\ldots,X_t]]/I$.  Put
$B'=\overline{k}[[X_1,\ldots,X_t]]/I\overline{k}[[X_1,\ldots,X_t]]$.
The map of power-series rings is flat, by the local flatness criterion
for the variables and the flat coefficient-field extension.  Its base
change to the quotient has the asserted residue field and maximal ideal.

In the remaining case the residue characteristic is a prime $p$.
Cohen structure supplies a presentation $B=V[[X_1,\ldots,X_t]]/I$,
where $V$ is a complete discrete valuation ring with uniformizer $p$
and residue field $k$.  We construct a discrete valuation extension
with unchanged uniformizer and residue field $\overline{k}$, without
assuming that $k$ is perfect.  Given a simple algebraic residue-field
extension, lift its monic irreducible polynomial to $f\in V[X]$.
The algebra $V[X]/(f)$ is finite free over $V$, its reduction modulo $p$
is the desired field, and all its maximal ideals lie over $(p)$.
It is therefore local with maximal ideal generated by $p$.
Since $p$ is a nonzerodivisor, it is a one-dimensional regular local
ring, hence a discrete valuation ring.  This argument also applies
to inseparable residue extensions.

Perform this construction along a tower of simple algebraic extensions
exhausting $\overline{k}$, taking unions at limit stages.  In every such
union, a nonzero element is a unit times a nonnegative integral power
of $p$.  Every nonzero ideal is generated by an element of least
valuation, so the union is still a discrete valuation ring.  Its
completion $V'$ has uniformizer $p$ and residue field $\overline{k}$.
The map $V\to V'$ is flat because $V'$ is torsion-free over $V$.
Now take
\[
 B'=V'[[X_1,\ldots,X_t]]/IV'[[X_1,\ldots,X_t]].
\]
The same local flatness criterion proves the required flatness.
This is the usual residue-extension construction; compare
\cite[Lemma~2.2]{Ma}.

For either construction, tensor a composition series of a finite-length
$B$-module with $B'$.  Flatness preserves exactness, and every simple
factor becomes $B'/\mathfrak b'$.  This proves the length assertions
and, applied to $V/\mathfrak b^aV$, the displayed identity.  Equality of
Hilbert--Samuel functions gives equality of module dimensions.  If $B$
is regular of dimension $d$, its $d$ maximal-ideal generators generate
$\mathfrak b'$, and $\dim B'=d$, so $B'$ is regular.  Finally flat base
change identifies the Tor modules after extension; their finite
lengths are preserved by the preceding argument.
\end{proof}

\subsection{Regular bases and the tower length}

For a perfect field $k$ of characteristic $p$, let $W(k)$ denote its
ring of Witt vectors.

\begin{lemma}[Regular normalization bases]\label{len:normalization}
Let $(D,\mathfrak d,k)$ be a complete Noetherian local domain of
dimension $h>0$, where $k$ is perfect of characteristic $p>0$.
There is a finite, generically separable inclusion $A\subseteq D$,
inducing the identity on residue fields, with
\[
 A=k[[t_1,\ldots,t_h]]\quad\text{if }\operatorname{char}D=p,
 \qquad
 A=W(k)[[t_2,\ldots,t_h]]\quad\text{if }\operatorname{char}D=0.
\]
In the second case set $t_1=p$.  In both cases $t_1,\ldots,t_h$
is a parameter system of $D$ and a regular parameter system of $A$.
\end{lemma}

\begin{proof}
The positive-characteristic assertion is the Cohen--Gabber
normalization theorem \cite{CohenGabber}.  In mixed characteristic,
choose a coefficient ring $W(k)\subseteq D$ and extend $p$ to a
parameter system $p,t_2,\ldots,t_h$.  Completeness and the parameter
condition make the resulting power-series map module finite.
Its kernel is zero: the source and target have dimension $h$, whereas
a nonzero ideal of the regular local domain
$W(k)[[t_2,\ldots,t_h]]$ lowers dimension.  The fraction-field
extension is separable in characteristic zero.
\end{proof}

Fix one of the bases in Lemma~\ref{len:normalization} and compatible
$p$-power roots of its parameters.  Write
\begin{equation}\label{len:tower-definition}
 A_n=A[t_1^{1/p^n},\ldots,t_h^{1/p^n}],\qquad
 \Lambda=\varinjlim_n A_n.
\end{equation}
The tower is uncompleted.  Only in mixed characteristic do we also
use its $p$-adic completion, denoted $A_\infty$.

\begin{lemma}[The tower and its length]\label{len:tower}
Each $A_n$ is a complete regular local ring with residue field $k$.
For $b\geq a$, the map $A_a\to A_b$ is finite free of rank
$p^{(b-a)h}$; in particular $A\to A_n$ has rank $p^{nh}$.
On the category of $\Lambda$-modules killed by some power of
$\mathfrak m_A$ there is a nonnegative additive function
$\lambda_\infty$ with the following properties.
\begin{enumerate}
\item A finitely presented module descends to a finite-length
      $A_n$-module $V_n$, and
      \begin{equation}\label{len:finite-stage}
       \lambda_\infty(V_n\otimes_{A_n}\Lambda)
                      =p^{-nh}\length_{A_n}(V_n).
      \end{equation}
\item For a finitely generated module $V$, its length is the infimum
      of the lengths of finitely presented modules in this category
      admitting a surjection onto $V$.
\item For an arbitrary module $V$ in this category,
      \begin{equation}\label{len:supremum}
       \lambda_\infty(V)=
       \sup\{\lambda_\infty(V_0):V_0\subseteq V
                                  \text{ is finitely generated}\}.
      \end{equation}
\end{enumerate}
The value on an arbitrary module may be infinite.  In mixed
characteristic, restriction from $A_\infty$ to $\Lambda$ gives the same
length on modules killed by a fixed power of $\mathfrak m_A$.
\end{lemma}

\begin{proof}
In characteristic $p$, the stages are power-series rings in the chosen
roots.  In mixed characteristic they are power-series rings over the
complete discrete valuation ring $W(k)[p^{1/p^n}]$.
The monomials in the successive roots, with each exponent less than
the relevant root degree, form bases for all transition extensions.
This proves the assertions about regularity, ranks, and residue fields.

Finite presentations over the union descend to a stage.  The additional
relations expressing annihilation by a fixed power of
$\mathfrak m_A$ also descend after increasing that stage, because
there are finitely many generators and relations to consider.  The
resulting finite-stage module has finite length.  Formula
\eqref{len:finite-stage} is unchanged on enlarging the stage, by
flatness and the computed free ranks.  The infimum and supremum
extensions, including their additivity and compatibility with this
formula, are Faltings' normalized-length formalism
\cite[Appendix~2]{Faltings}, in the formulation of
\cite[Definition~2.2.1 and Propositions~2.2.3, 2.2.6(a)]{CLMST}.
The same residue field implies that $A_n$-length and $A$-length
agree on finite-length $A_n$-modules.

In mixed characteristic $\Lambda$ is $p$-torsion-free, and the natural
map induces $\Lambda/p^a\Lambda\simeq A_\infty/p^aA_\infty$ for every
$a\geq1$.  A module killed by $\mathfrak m_A^a$ is killed by $p^a$.
Its submodules and finite generating sets are therefore identical
over the two bases.  In computing the infimum over finitely presented
covers of such a module, one may replace a cover $V$ by
$V/\mathfrak m_A^aV$; this is still finitely presented.  These covers
also correspond over the identical quotient rings, proving the last
assertion.
\end{proof}

\begin{lemma}[Zero-length limits]\label{len:zero-length-limits}
Suppose all modules under consideration are killed by one fixed power
of $\mathfrak m_A$.  Sums, direct sums, and directed colimits of
zero-length modules have length zero.  For an increasing union,
\[
 \lambda_\infty\Bigl(\bigcup_a V_a\Bigr)
                         =\sup_a\lambda_\infty(V_a).
\]
If $(V_a)_{a\geq1}$ is a direct system with
$\lambda_\infty(V_a)\to0$, then
$\lambda_\infty(\varinjlim_a V_a)=0$, without any injectivity
assumption on the transition maps.
\end{lemma}

\begin{proof}
A finite set of elements of a direct sum is contained in a finite
subsum.  Additivity and \eqref{len:supremum} give the assertion for
direct sums; sums and directed colimits are quotients of direct sums.
For an increasing union, every finite generating set lies in one
stage.  Finally a finitely generated submodule of $\varinjlim_a V_a$
is contained in the image of every sufficiently late $V_a$, by lifting
its generators to a common stage.  Monotonicity bounds its length by
$\lambda_\infty(V_a)$ at all such stages.  Apply
\eqref{len:supremum}.  The common annihilator ensures throughout
that these modules belong to the length category.
\end{proof}

\subsection{Length algebras and parameter homology}

The tower length will now measure modules over an algebra of the given
domain.  Its parameter identities have two uses: the base parameters
control the Euler comparisons, while arbitrary parameter systems give
the positive quotients used in Section~\ref{sec:positivity}.
For an ideal $I$ primary to the maximal ideal of a local ring $D$, write
$e(I,D)$ for its Hilbert--Samuel multiplicity.

Related normalized-length constructions, in which parameter-quotient
lengths recover multiplicities and positive parameter Koszul homology has
normalized length zero, appear in \cite[Theorem~3.1]{OpenAILech2026}.
Here we retain the chosen base $A\subseteq D$ and the finite-stage,
infimum, and supremum properties on underlying tower modules.

\begin{theorem}[Normalized-length input]\label{len:input}
Let $(D,\mathfrak d,k)$ and a chosen finite generically separable base
$A\subseteq D$ be as in Lemma~\ref{len:normalization}, and let
$\Lambda$ be the tower \eqref{len:tower-definition}.  There is a
$D$-algebra $C_D$, carrying a compatible $\Lambda$-action, such that
restriction of scalars defines an additive nonnegative length
$\lambda_D=\lambda_\infty$ on its modules annihilated by a power of
$\mathfrak d$.  For every parameter system $z_1,\ldots,z_h$ of $D$
and every integer $v\geq1$,
\begin{align}
 \lambda_D H_i(z_1^v,\ldots,z_h^v;C_D)&=0\qquad(i>0),
                                      \label{len:koszul-vanishing}\\
 \lambda_D\bigl(C_D/(z_1^v,\ldots,z_h^v)C_D\bigr)
                    &=v^h e((z_1,\ldots,z_h),D).
                                      \label{len:parameter-length}
\end{align}
The right side is positive and finite.  The algebra and length are
independent of the parameter system $z$.
For the regular parameters of $A$, in particular,
\begin{equation}\label{len:base-parameter-length}
 \lambda_D\bigl(C_D/(t_1^v,\ldots,t_h^v)C_D\bigr)
                         =v^h\rank_A D.
\end{equation}
The finite-stage, infimum, and supremum properties of
Lemma~\ref{len:tower} remain available on underlying tower modules;
they do not require that the relevant submodules or covers be
$C_D$-modules.
\end{theorem}

\begin{proof}
We give the positive-characteristic construction and then specify the
mixed-characteristic literature inputs.

Assume first that $D$ has characteristic $p$, and put
$C_D=D_{\mathrm{perf}}=\bigcup_n D^{1/p^n}$ inside a fixed algebraic
closure of its fraction field. Here $\Lambda=A_{\mathrm{perf}}$.
Write $K=\operatorname{Frac}(A)$, $L=\operatorname{Frac}(D)$, and
$K'=K_{\mathrm{perf}}$.  For $q=p^n$ put
\[
 B_q=\Lambda\otimes_{A^{1/q}}D^{1/q}.
\]
Flatness of $\Lambda$ over $A^{1/q}$ and torsion-freeness of
$D^{1/q}$ embed this tensor product into
$K'\otimes_{K^{1/q}}L^{1/q}$.  Generic separability makes the latter
a field, by linear disjointness from a purely inseparable extension.
Thus $B_q=\Lambda D^{1/q}\subseteq C_D$, and these subrings increase
to $C_D$.

Let $H$ be any $\mathfrak d$-primary ideal and set
$V_q=B_q/HB_q$ and $W_q=\operatorname{im}(V_q\to C_D/HC_D)$.
Write $H^{[q]}=(a^q:a\in H)$ for its Frobenius power.
Finite-stage normalization and the $q$-th-power isomorphism give
\[
 \lambda_\infty(V_q)=q^{-h}\length_D(D/H^{[q]}).
\]
Although $B_q\subseteq C_D$, the map on these quotients need not be
injective.  To recover the length of $C_D/HC_D$ from the displayed
formula, we control the kernels of $V_q\to W_q$ by a single trace
denominator.  There is a nonzero $g\in A$ with
\begin{equation}\label{len:conductor}
 gC_D\subseteq B_1,\qquad g^{1/q}C_D\subseteq B_q.
\end{equation}
Indeed, choose a $K$-basis $b_1,\ldots,b_r$ of $L$ in $D$ and its
trace-dual basis $\beta_1,\ldots,\beta_r$.  A common denominator
$g\in A\setminus\{0\}$ makes every $g\beta_j$ lie in $D$.
The field $LK'$ is perfect and equals $L_{\mathrm{perf}}$.
The ring $\Lambda$ is normal, since an integral equation and the
fractions occurring in it descend to a normal stage $A^{1/q}$.
For $c\in C_D$, each
$\operatorname{Tr}_{LK'/K'}(g\beta_jc)$ is integral over $\Lambda$
and belongs to $K'$, hence belongs to $\Lambda$.  The trace-dual
expansion of $gc$ proves the first inclusion in
\eqref{len:conductor}.  Taking $q$-th roots proves the second.

The kernel of $V_q\to W_q$ is killed by $g^{1/q}$, by
\eqref{len:conductor}.  Kernel and cokernel of an endomorphism of a
finite-length object have equal length.  Applied to multiplication
by $g^{1/q}$, this yields
\[
 \lambda_\infty\ker(V_q\to W_q)
 \leq q^{-h}\length_D D/(H^{[q]},g)=O(q^{-1}).
\]
For the estimate, if $H$ has $s$ generators then
$H^{s(q-1)+1}\subseteq H^{[q]}$.  Thus $H^{[q]}$ contains a power
of $\mathfrak d$ with exponent $O(q)$, whereas
$\dim D/(g)\leq h-1$ if $g$ is a nonunit.  Hilbert--Samuel growth
gives the asserted bound; for a unit $g$ the quotient is zero.
The $W_q$ increase to $C_D/HC_D$ and are uniformly primary-torsion.
The Hilbert--Kunz limit and Lemma~\ref{len:zero-length-limits} show
\begin{equation}\label{len:perfection-quotient}
 \lambda_D(C_D/HC_D)=e_{\mathrm{HK}}(H,D).
\end{equation}
Here we use Monsky's existence theorem for Hilbert--Kunz multiplicity
\cite{Monsky}; see also \cite[Section~2.2]{Ma}.

We have identified primary-quotient lengths with Hilbert--Kunz
multiplicities.  For a parameter ideal, it remains to obtain the
Hilbert--Samuel value and zero positive Koszul homology.  If $z$ is
a parameter system of $D$, finite-stage flatness gives
\[
 \lambda_\infty H_i(z;B_q)
          =q^{-h}\length_D H_i(z_1^q,\ldots,z_h^q;D).
\]
The Koszul complex $K(z;D)$ has finite free terms in degrees
$0,\ldots,h$ and finite-length homology.  The Frobenius homology
theorem of Roberts, in the form \cite[Theorem~2.4]{Ma}, says that
the displayed quantity tends to zero for $i>0$; the original result
is \cite{RobertsIntersectionTheorems}. Its hypotheses
require precisely a complete characteristic-$p$ ring and a
finite-free complex of length its dimension with finite-length
homology.  Filtered colimits commute with homology, and these
Koszul homology modules have the common annihilator $(z)$.
Lemma~\ref{len:zero-length-limits} proves
\eqref{len:koszul-vanishing}, also with $z$ replaced by $z^v$.
Finally the Koszul Euler characteristic formula gives
\[
 \sum_{i=0}^h(-1)^i\length_D H_i(z_1^q,\ldots,z_h^q;D)
                           =q^h e((z),D).
\]
After division by $q^h$, the positive-degree terms tend to zero.
Together with \eqref{len:perfection-quotient} and the multiplicity
power rule this proves \eqref{len:parameter-length}.

Now suppose $D$ has mixed characteristic.  Let $A_\infty$ be the
$p$-adic completion of $\Lambda$ and define
\[
 C_D=(D\otimes_A A_\infty)_{\mathrm{perfd}},
\]
the perfectoidization of \cite[Notation~4.0.1]{CLMST}, built on
Bhatt--Scholze's construction \cite[Proposition~8.5]{BhattScholze}.
The stated hypotheses there are exactly that $D$ is a complete
mixed-characteristic local domain with perfect residue field and
the chosen finite Witt power-series base $A$.
Choose $g\in A\setminus\{0\}$ so that $A[1/g]\to D[1/g]$ is
finite etale.  Corollary~4.0.4 of \cite{CLMST} annihilates positive
Koszul homology for every parameter system of $A_n\otimes_A D$
by the almost ideal $(g)_{\mathrm{perfd}}$.
At $n=0$ this includes every parameter system of $D$.
Proposition~4.0.10 of that paper makes these uniformly
$\mathfrak m_A$-power-torsion homology modules have normalized
length zero.  Proposition~4.0.14, again at $n=0$, identifies the
length of the parameter quotient with its Hilbert--Samuel
multiplicity.  These two results, applied also to $z^v$, prove
\eqref{len:koszul-vanishing} and \eqref{len:parameter-length}
for this single algebra $C_D$.  Lemma~\ref{len:tower} permits
restriction to the uncompleted tower throughout.

In either characteristic, a module killed by a power of
$\mathfrak d$ is killed by a power of $\mathfrak m_A$, so all
claimed lengths are defined.  Additivity and the tower properties
are those already established.  Lastly the multiplicity rank formula,
and equality of the residue fields, give
$e((t),D)=\rank_A D\,e((t),A)=\rank_A D$.
This proves \eqref{len:base-parameter-length}.
\end{proof}

The base-parameter case extends to every finite-length module over
$A$.  The same parameter-quotient positivity also prevents the residue
quotient of $C_D$ from having length zero; this will supply the positive
residue-length factor in the final degree-zero estimate.

\begin{corollary}[Finite-length base change]\label{len:finite-length-base-change}
In Theorem~\ref{len:input}, put $r=\rank_A D$.
For every finite-length $A$-module $V$,
\[
 \lambda_D\Tor_i^A(V,C_D)=0\quad(i>0),\qquad
 \lambda_D(V\otimes_A C_D)=r\length_A(V).
\]
Moreover $\lambda_D(C_D/\mathfrak d C_D)>0$.
\end{corollary}

\begin{proof}
For $V=k$, the first two assertions are
\eqref{len:koszul-vanishing} and \eqref{len:base-parameter-length},
using the Koszul resolution over regular $A$.
Induct on a composition series of $V$.  The Tor long exact
sequences and length additivity discard all positive-degree terms
and show that degree-zero length increases by $r$ at each step.
All modules here have bounded primary torsion.

For the last assertion choose a parameter ideal $H$ of $D$.
Tensor a composition series of $D/H$ with $C_D$.
Right exactness gives a filtration of $C_D/HC_D$ whose factors
are quotients of $C_D/\mathfrak dC_D$.  Therefore
\[
 0<e(H,D)=\lambda_D(C_D/HC_D)
       \leq\length_D(D/H)\lambda_D(C_D/\mathfrak dC_D).
\]
This proves the required strict positivity.
\end{proof}
